% !TeX root = ../main.tex
\section{Индивидуальные значения и результаты начального моделирования}
Изменим параметры сигнала и гетеродина в соответсвии с номером по журналу ($N_{\text{пж}} = 4$) с учётом отстройки частоты входного монохроматического сигнала на $0{,}1$~МГц (Таблица~\ref{tab:initial_data}).

\begin{table}[H]
    \small
    \caption{Параметры устройства цифровой демодуляции (вариант 4)}
    \label{tab:initial_data}
    \begin{tabularx}{\textwidth}{|X|c|c|}
        \hline
        \textbf{Наименование параметра} & \textbf{Формула} & \textbf{Значение} \\ \hline
        Промежуточная частота входного сигнала, $f_C$ & $130 + N_{\text{пж}} \cdot 10 + 0{,}1$ & $170{,}1$ МГц \\ \hline
        Частота дискретизации АЦП, $f_S$ & $100 + N_{\text{пж}} \cdot 10$ & $140{,}0$ МГц \\ \hline
        Количество бит АЦП, $n_a$ & $5 + N_{\text{пж}}$ & 9 бит \\ \hline
        Центральная частота цифрового гетеродина (ЦГ), $f_G$ & $130 + N_{\text{пж}} \cdot 10$ & $170{,}0$ МГц \\ \hline
        Количество бит цифрового гетеродина, $n_g$ & $5 + N_{\text{пж}}$ & 9 бит \\ \hline
    \end{tabularx}
\end{table}

Проведем моделирования сигнала при отношениии ОСШ равным 75 дБ и при фазовой ошибки гетеродина равной $2^\circ$

\begin{figure}[H]
    \centering
    \includegraphics[width=0.8\textwidth]{2_1_i.png}
    \caption{Входной сигнал перед АЦП}
\end{figure}

\begin{figure}[H]
    \centering
    \includegraphics[width=0.8\textwidth]{2_2_i.png}
    \caption{Спектр входного сигнала}
\end{figure}

Поскольку частота дискретизации $f_{s}=140 МГц$ меньше удвоенной несущей частоты, намеренно нарушается условие теоремы Котельникова. В результате спектр полезного сигнала переносится на разностную частоту |$170.1 - 140$|$ =30.1 МГц $. Также можно заметить, что шум, который мы ввели вместо сигналом равномерно распределен. Рассмотрим  следующие этапы преобразований сигнала - после АЦП.

\begin{figure}[H]
    \centering
    \includegraphics[width=0.8\textwidth]{2_3.png}
    \caption{Сигнал после АЦП}
\end{figure}


\begin{figure}[H]
    \centering
    \includegraphics[width=0.8\textwidth]{2_4_i.png}
    \caption{Спектр после АЦП}
\end{figure}

Как мы можем заметить, существуют ошибки округления в АЦП, а также перекрытие спектра сигнала с самим собой из-за невыполнения требований теоремы Котельникова. Рассмотрим  следующие этапы преобразований сигнала - сигнал цифрового генератора с небольшими расстройкой по фазе (в квадратурном канале присутсвует рассогласование $+2^\circ$ по фазе ).

\begin{figure}[H]
    \centering
    \includegraphics[width=0.8\textwidth]{2_5.png}
    \caption{Сигнал Цифрового Гетеродина}
\end{figure}


\begin{figure}[H]
    \centering
    \includegraphics[width=0.8\textwidth]{2_6_i.png}
    \caption{Спектр Цифрового Гетеродина}
\end{figure}

На рисунках относящихся к цифровому гетеродину с фазавым дисбалансом между каналами можно заметить, что у генератора присутствует дополнительные гармоники (шпоры), что негативно сказывается над демодуляцией сигнала показанной ниже. Эти шпоры носят периодический характер, из-за ошибок округления, генераций нечетных гармоник, эффектом наложения, приведенные ниже

\begin{itemize}
        \item 3-я гармоника ($3 \times 30 = 90$ МГц) переносится на $90 - 140 = -50$  МГц.
        \item 5-я гармоника ($5 \times 30 = 150$ МГц) переносится на $150 - 140 = +10$ МГц.
        \item 9-я гармоника ($9 \times 30 = 270$ МГц) переносится на $270 - 2 \times 140 = -10$ МГц.
        \item 11-я гармоника ($11 \times 30 = 330$ МГц) переносится на $330 - 2 \times 140 = +50$ МГц.
    \end{itemize}

\begin{figure}[H]
    \centering
    \includegraphics[width=0.8\textwidth]{2_7.png}
    \caption{Осцилограмма каналов IQ в DEM}
\end{figure}


\begin{figure}[H]
    \centering
    \includegraphics[width=0.8\textwidth]{2_8_i.png}
    \caption{Спектр сигнала на выходе DEM}
\end{figure}

Представленные рисунки наглядно показывают результат перемножения  сигналов оцифровки и цифрового гетеродина, на нижние частоты, но из-за наложения спектров мы видим и другие копии сигналов.